% !TEX root = ../Thesis.tex

\section{What Holds}

Six claims survive across the thesis, each with its own supporting result rather than a
shared assumption.

\vspace{10pt}

The trial's average effect is confirmed. Five estimator families that share no code path
beyond preprocessing all recover a bleeding benefit of abbreviated therapy of
$4$--$7$ percentage points and an ischemic effect statistically indistinguishable from
zero, matching the published trial (Chapter 6).

\vspace{10pt}

No exploitable heterogeneity emerges from the direct CATE ranking. 
The validation compares the dispersion of the
estimated CATE, and the doubly-robust ranking built from it, against a floor obtained by
permuting the treatment label rather than against zero, so a positive dispersion alone is
not treated as evidence of heterogeneity. Under that comparison none of the five
estimators separates from noise on any endpoint (Chapter 6).

\vspace{10pt}

The same conclusion is not reached a second time ithout estimating a single CATE
Chapter~\ref{ch:guided-enrollment-applied}'s guided-enrollment mechanisms move the
estimated-CATE plane exactly as designed, and
that movement is confirmed by a reference model that took no part in selection, yet the
observed bleeding and ischemic event rates of the resulting cohorts do not differ
reliably from a matched random cohort at either exploration setting tested, and the two
arms that do reach nominal significance move in the direction opposite to what a
trade-off-seeking mechanism would need.

\vspace{10pt}

The limit is in the cohort, not in the tools, and this is verified rather than assumed.
The positive control shows the estimators recover a real average effect when one is
present; the minimum detectable $\delta$ of Chapter 6 shows the same protocol recovers
\emph{synthetic} heterogeneity of known magnitude once it is injected. Both checks would
have caught a broken pipeline, and neither does.

\vspace{10pt}

The mechanism behind the negative result is known. MASTER DAPT's eligibility criterion is
high bleeding risk, which compresses the bleeding-risk axis before any model sees a
patient; this is why bleeding, with 506 events, is predicted worse (AUROC 0.62) than
cardiovascular death, with 81 events (AUROC 0.76), in Chapter 5, and it is why the causal
forest has little covariate-space variation left to exploit on that axis in Chapter 6.
The aggregate ischemic score compounds the problem on the other axis: MACCE's near-zero
average effect hides opposite-signed MI and stroke components rather than a genuine
absence of effect on either.

\vspace{10pt}

Two methodological results hold independently of the case study. Selecting which patients
enter a trial by their estimated score preserves randomization, because the score depends
only on covariates measured before assignment; allocating patients to an arm by the same
score does not, and needs the tempered, rejection-sampled rule of Chapter 3 to stay
identified. And a linear combination of two effects of unequal variance collapses toward
the larger one: the win--win score of Chapter~\ref{ch:guided-enrollment-applied} correlates $\rho=+0.91$ with the bleeding
CATE alone, and its diagonal-geometry counterpart produces an identical ranking whether it
is built from the win--win or the trade-off scalar, because both read the same two
underlying numbers.

\section{What Does Not Hold}

The minimum detectable $\delta$ on bleeding is $0.05$ on the absolute-risk scale, as large
as the trial's entire average effect. What these data exclude is heterogeneity of that
magnitude, not heterogeneity in general: a modifier half as large, still clinically
relevant, would have been missed 20\% of the time. The correct claim throughout this
thesis is \emph{not found}, never \emph{does not exist}.

\vspace{10pt}

Three of the five endpoints cannot be assessed at all. MI, cardiovascular death and stroke
carry between 35 and 109 events, and the sensitivity protocol shows that detecting effect
modification on MI or cardiovascular death would require an odds ratio between 3 and 5,
a magnitude nobody expects in this setting. On these three endpoints the honest statement
is \emph{not measurable on this cohort}, and the stroke row of the sensitivity analysis is
not yet interpretable at all: under the null it beats its own noise floor in 85\% of
seeds instead of the expected 50\%, because the simulated event count at that endpoint's
rate is too low to support the protocol.

\vspace{10pt}

The conflict-subgroup profile of Chapter 6, 687 patients whose observed bleeding
difference under the two arms reaches 20.8 percentage points, is an in-sample,
descriptive result: the score that selected those patients was fitted on the same
patients and the same outcomes, and none of them is distinguishable from the cohort
average once a bootstrap confidence interval is attached.

\vspace{10pt}

Chapter~\ref{ch:guided-enrollment-applied}'s own scope is narrower than its headline result suggests. The main comparison
is read at a checkpoint where 92\% of the cohort is already enrolled, a regime that leaves
little room for any policy to diverge from the population it is drawn from, which biases
the comparison toward the null reported rather than away from it. The one mirror-test
regression that does show the expected sign, in the tempered configuration, loses both its
sign and its significance in the greedy one, on only five points; it does not survive as a
stable result. And no matched-size, non-adaptive randomized-arrival baseline exists for
this cohort, so Chapter~\ref{ch:guided-enrollment-applied} can say that guided enrollment does not change which real
outcomes occur at a given enrolled size, but it cannot say whether any mechanism would
reach a given precision with fewer patients than random enrollment, which was the second
half of the question Chapter 3 originally posed.

\vspace{10pt}

Beyond the specific analyses, the thesis carries three general limitations. There is no
external validation on a cohort independent of MASTER DAPT. Multiple testing across five
endpoints and five estimators is declared throughout but not formally corrected. The
computational environment is not fully packaged for independent reproduction.

\section{To Which Population the Same Question Should Be Put}

The negative result is useful only if it points somewhere. If the cause is range
restriction on bleeding risk, the same question needs a cohort where that axis has not
already been compressed by the eligibility criterion, which excludes a trial that, like
MASTER DAPT, enrolls exclusively high-bleeding-risk patients. The ischemic side needs
power rather than a wider axis: 35 strokes and 81 cardiovascular deaths cap what any
estimator can find regardless of its sophistication, so a longer follow-up or a
population selected for ischemic rather than bleeding risk is the more direct fix.

\vspace{10pt}

A trial designed to answer this question would also need to be sized differently. MASTER
DAPT, like most trials, is powered for the average effect; the minimum detectable
$\delta$ of Chapter 6 is precisely the tool that turns \emph{how much heterogeneity would
this design see} into a number that can be checked before a single patient is enrolled,
rather than discovered afterward as this thesis had to. It is on such a cohort, not on
this one, that the guided-enrollment machinery of Chapter 3 and
Chapter~\ref{ch:guided-enrollment-applied} would have something to
find: the mechanisms, the reference-model discipline, and the two documented failure
modes, allocation without tempering and linear scores collapsing toward the higher-variance
axis, are all already built and already measured.

\section{Further Work}

The minimum detectable $\delta$ needs closing out before it can be cited without caveats:
the baseline risk model $\mu(x)$ underestimates bleeding prevalence, which makes the
reported $\delta^*$ conservative in an unquantified direction, and the stroke row needs a
protocol that does not degenerate at that event count. Extending the sensitivity analysis
to the other four estimators, not only the causal forest, would show whether $\delta^*$ is
a property of the data or of one splitting rule.

\vspace{10pt}

Formal correction for multiple testing across the five endpoints and five estimators is
still open, as is a cross-fitted AIPW estimate of the trial-level ATE with confidence
intervals compared directly against the published trial's own.

\vspace{10pt}

Two extensions follow directly from Chapter~\ref{ch:guided-enrollment-applied}'s own limitations. The whole-pool
comparison should be re-run at an earlier checkpoint, before the cohort is 92\% enrolled,
so that any real divergence between policies is not structurally suppressed; and a
matched-size, non-adaptive randomized-arrival baseline should be built so that precision,
not only cohort composition, can be compared at equal enrolled $n$ — the comparison
Chapter 3 originally set out to make and Chapter~\ref{ch:guided-enrollment-applied} could not complete.

\vspace{10pt}

The most direct extension remains replicating the whole protocol on a cohort not selected
for bleeding risk, since Chapter 6's explanation of the negative result predicts that this
is where the same machinery would stop finding nothing.
